\section*{الفصل \ref{ch:b1:nernst} --- توازنات الأكسدة والاختزال ومعادلة نرنست}

\begin{solution}{exo:b1:nernst:1}
N في \ce{NH3}: \babelsublr{$-$III}؛ وفي \ce{NO3-}: \babelsublr{$+$V}. S في \ce{SO4^2-}: \babelsublr{$+$VI}. O في
\ce{H2O2}: \babelsublr{$-$I}. Cr في \ce{Cr2O7^2-}: \babelsublr{$+$VI}. Cl في \ce{ClO-}: \babelsublr{$+$I}. Fe في
\ce{Fe3O4}: $+8/3$ في المتوسط (حديد(II) واحد وحديدان(III)). C في
\ce{CH3OH}: $x + 4(+1) + (-2) = 0$، \babelsublr{$-$II}.
\end{solution}

\begin{solution}{exo:b1:nernst:2}
\ce{Cr2O7^2- + 14H+ + 6e- -> 2Cr^3+ + 7H2O}؛
\ce{H2O2 + 2H+ + 2e- -> 2H2O}؛
\ce{O2 + 2H+ + 2e- -> H2O2}؛
\ce{SO4^2- + 4H+ + 2e- -> SO2 + 2H2O}.
\end{solution}

\begin{solution}{exo:b1:nernst:3}
$E = -0.76 + 0.030\log[\ce{Zn^2+}]$؛ $E = 1.36 + 0.030\log(p(\ce{Cl2})/[\ce{Cl-}]^2)$؛
$E = 1.23 + 0.015\log(p(\ce{O2})[\ce{H+}]^4)$؛ $E = 0.27 - 0.059\log[\ce{Cl-}]$
(الضغوط بالبار). وفي كبريتات الزنك ذات التركيز \qty{0.010}{mol/L}:
$E = -0.76 + 0.030 \times (-2) = -0.82$~V.
\end{solution}

\begin{solution}{exo:b1:nernst:4}
$E(\ce{Zn}) = -0.76 + 0.030\log 0.10 = -0.79$~V، $E(\ce{Cu}) = 0.34 +
0.030\log 0.010 = 0.28$~V. الجهد $0.28 - (-0.79) = 1.07$~V. والزنك، ذو
الجهد الأدنى، هو القطب السالب، حيث تجري الأكسدة: فهو
\omterm{def:b1:nernst:cell}{المصعد}.
\end{solution}

\begin{solution}{exo:b1:nernst:5}
\ce{MnO4- + 5Fe^2+ + 8H+ -> Mn^2+ + 5Fe^3+ + 4H2O}، $n = 5$،
$\log K = 5(1.51 - 0.77)/0.059 = 63$: تفاعل كمّي وسريع، ومعايرة
كلاسيكية (البرمنغنات كاشف نفسها).
\end{solution}

\begin{solution}{exo:b1:nernst:6}
تتفاعل الفلزات التي $E^\circ$ لها دون 0: Mg، Zn، Fe، Ni، Pb (الرصاص
ببطء فقط: الفرق صغير)؛ ولا يتفاعل Cu وAg. \ce{Mg + 2H+ -> Mg^2+ + H2}، $\log K = 2 \times 2.36/0.059 = 80$؛
\ce{Fe + 2H+ -> Fe^2+ + H2}، $\log K = 2 \times 0.41/0.059 = 13.9$.
\end{solution}

\begin{solution}{exo:b1:nernst:7}
$n = 2$، $\log K = 2(0.53 - 0.02)/0.059 = 17$: كمّي.
\end{solution}

\begin{solution}{exo:b1:nernst:8}
$3E^\circ(\ce{Fe^3+}/\ce{Fe}) = 2(-0.41) + 0.77$، $E^\circ = -0.02$~V.
\ce{2Fe^3+ + Fe -> 3Fe^2+}: \omterm{def:b1:nernst:couple}{المؤكسد} \ce{Fe^3+} (0.77) فوق \omterm{def:b1:nernst:couple}{المختزل}
\ce{Fe} ($-0.41$)، $n = 2$، $\log K = 2(0.77 + 0.41)/0.059 = 40$. فبرادة
الحديد تمنع محلول الحديد(II) من التحول إلى حديد(III).
\end{solution}

\begin{solution}{exo:b1:nernst:9}
\ce{O2 + 4H+ + 4e- -> 2H2O}: $E = 1.23 + \frac{0.059}{4}\log [\ce{H+}]^4 =
1.23 - 0.059\,\mathrm{pH}$. \ce{2H+ + 2e- -> H2}: $E = \frac{0.059}{2}\log
[\ce{H+}]^2 = -0.059\,\mathrm{pH}$. والفرق $1.23$~V عند كل pH:
\omterm{def:b1:nernst:cell}{فجهد خلية} الوقود لا يتعلق بقيمة pH.
\end{solution}

\begin{solution}{exo:b1:nernst:10}
$E^\circ(\ce{Cu+}/\ce{Cu}) = 0.52 > E^\circ(\ce{Cu^2+}/\ce{Cu+}) = 0.16$:
\omterm{def:b1:nernst:couple}{فالمؤكسد} \ce{Cu+} في المزدوجة العليا يتفاعل مع \omterm{def:b1:nernst:couple}{المختزل} \ce{Cu+} في
السفلى. $\log K = (0.52 - 0.16)/0.059 = 6.1$، $K =
[\ce{Cu^2+}]/[\ce{Cu+}]^2$، فيكون $[\ce{Cu+}] = \sqrt{0.10/10^{6.1}} =
\qty{2.8e-4}{mol/L}$.
\end{solution}

\begin{solution}{exo:b1:nernst:11}
$U = 0.059\log(0.10/\num{1.0e-3}) = 0.118$~V. والقطب الموجب هو
الفضة في المحلول المركز، حيث يُختزل \ce{Ag+}؛ وفي المحلول
المخفف تتأكسد الفضة. فيصير المحلول المركز
أكثر تخفيفًا والمخفف أكثر تركيزًا؛ ويهبط الجهد
ويبلغ الصفر حين يتساوى التركيزان.
\end{solution}

\begin{solution}{exo:b1:nernst:12}
\ce{H2O2} هو \omterm{def:b1:nernst:couple}{مؤكسد} $\ce{H2O2}/\ce{H2O}$ (1.76~V) \omterm{def:b1:nernst:couple}{ومختزل}
$\ce{O2}/\ce{H2O2}$ (0.69~V): فهو يتفاعل مع نفسه، \ce{2H2O2 -> 2H2O + O2}،
$\log K = 2(1.76 - 0.69)/0.059 = 36$. والتفاعل مواتٍ لكنه بطيء جدًا
دون \omterm{def:b1:elementary-steps:catalyst}{حفاز} (\cref{ch:b1:elementary-steps})؛ والقوارير النظيفة، والحفظ
في البرد والظلام، تُبقيه بطيئًا.
\end{solution}

\begin{solution}{pb:b1:nernst:1}
\textbf{1.} Ag: 0، \babelsublr{$+$I}، \babelsublr{$+$I}؛ Cl: \babelsublr{$-$I} في الاثنين.
\textbf{2.} \ce{AgCl(s) + e- -> Ag(s) + Cl-}.
\textbf{3.} $E = E^\circ(\ce{AgCl}/\ce{Ag}) - 0.059\log[\ce{Cl-}]$.
\textbf{4.} \omterm{def:b1:extent-q-and-k:activity}{نشاطية} الجسمين الصلبين 1؛ ولا يوجد في المحلول إلا أيون
الكلوريد.
\textbf{5.} $E^\circ = 0.22$~V.
\textbf{6.} $0.22 + 0.118 = 0.34$~V.
\textbf{7.} \babelsublr{\ce{Pt} $|$ \ce{H2} (\qty{1}{bar}) $|$ \ce{H+}
$\|$ \ce{Cl-} (\qty{0.10}{mol/L}) $|$ \ce{AgCl} $|$ \ce{Ag}}، \omterm{def:b1:extent-q-and-k:activity}{ونشاطية} أيونات الهيدروجين 1.
\textbf{8.} قطب الفضة: $0.22 + 0.059 = 0.28$~\babelsublr{V $> 0$}.
\textbf{9.} $0.28$~V.
\textbf{10.} \ce{2AgCl + H2 -> 2Ag + 2H+ + 2Cl-}.
\textbf{11.} $\log K = 2 \times 0.22/0.059 = 7.5$.
\textbf{12.} من البلاتين (القطب السالب، حيث يتأكسد \ce{H2})
إلى الفضة.
\textbf{13.} $0.27 - 0.059\log 1.0 = 0.27$~V.
\textbf{14.} $E(\ce{Ag}) = 0.22 + 0.177 = 0.40$~\babelsublr{V $> 0.27$}: قطب الفضة
موجب، $U = 0.13$~V.
\textbf{15.} $U = 0.22 - 0.059\log[\ce{Cl-}] - 0.27 = -0.05 -
0.059\log[\ce{Cl-}]$ (بالفولت).
\textbf{16.} $\log[\ce{Cl-}] = -(0.115 + 0.05)/0.059 = -2.80$،
$[\ce{Cl-}] = \qty{1.6e-3}{mol/L}$، أي \qty{57}{mg/L}.
\textbf{17.} $\Delta c/c = \ln 10 \times 0.001/0.059 = 0.039$: نحو
\qty{4}{\percent} لكل ميليفولت.
\textbf{18.} حين يصير كلوريد العيّنة مماثلًا لما يحرره
كلوريد الفضة نفسه، $\sqrt{K_s} = \qty{1.3e-5}{mol/L}$؛
ويُستعمل القطب فوق ذلك بكثير، ابتداءً من نحو \qty{1e-4}{mol/L}.
\textbf{19.} $E = 0.80 + 0.059\log[\ce{Ag+}]$.
\textbf{20.} $[\ce{Ag+}] = K_s/[\ce{Cl-}]$.
\textbf{21.} $E = 0.80 + 0.059\log K_s - 0.059\log[\ce{Cl-}] = 0.80 -
0.059\,\mathrm{p}K_s - 0.059\log[\ce{Cl-}]$.
\textbf{22.} $E^\circ(\ce{AgCl}/\ce{Ag}) = E^\circ(\ce{Ag+}/\ce{Ag}) -
0.059\,\mathrm{p}K_s$ (\cref{prop:b1:nernst:second-kind}).
\textbf{23.} $0.80 - 0.059 \times 9.75 = 0.225$~V، وهو متفق مع
الجدول (0.22~V)، المحسوب على نحو مستقل من طاقات غيبس.
\textbf{24.} $E^\circ(\ce{AgCl}/\ce{Ag}) = $ \textbf{0.22~V}.
\end{solution}
